\documentclass[11pt]{article}
\usepackage[english]{babel}
\usepackage[utf8]{inputenc}
\usepackage{amsmath,amssymb}
\usepackage{graphicx}
\usepackage[colorinlistoftodos]{todonotes}
\usepackage[]{cancel}
\usepackage[margin=1in]{geometry}
\usepackage[affil-it]{authblk}
\usepackage[]{lmodern}
\usepackage[]{nicefrac}
\usepackage{tikz}
\usetikzlibrary{arrows, shapes.geometric, positioning, patterns, shapes.arrows}
\usepackage[]{hyperref}
\hypersetup{colorlinks = true, urlcolor = gray}

\usepackage{xparse}
\usepackage{amsmath}

\newcommand{\bxi}{\boldsymbol{\xi}}
\newcommand{\dt}{\, {\rm d}t}
\newcommand{\dX}{\, {\rm d}\mathbf{X}}
\newcommand{\dx}{\, {\rm d}\mathbf{x}}
\newcommand{\dVblank}{\, {\rm d}V}
\newcommand{\dxi}{\, {\rm d}\bxi}
\newcommand{\deta}{\, {\rm d}\bs\eta}
\newcommand{\dv}{\,dv}
\newcommand{\fr}{\underline{\nabla}}
\newcommand{\eps}{\varepsilon}

\NewDocumentCommand \snsI{ m m o }{%
    \IfNoValueTF{#3}{\underline{\mbf{#1}}_{#2}}%
    {
      \underline{\mbf{#1}}_{#2}(#3)%
    }
}

\NewDocumentCommand \stateOneN{ m }{
    \underline{#1}%
}

\NewDocumentCommand \stateTwoN{ m o }{%
    \IfNoValueTF{#2}{\stateOneN{#1}}
    {
    \stateOneN{#1}\an{#2}%
    }
}

\NewDocumentCommand \sN{ m o o }{%
    \IfNoValueTF{#3}{\stateTwoN{#1}[#2]}%
    {
    \underline{#1}(#3)\an{#2}%
    }
}

\NewDocumentCommand \dV{ o }{%
    \IfNoValueTF{#1}{\dVblank}%
    {
        \dV_{#1}
    }
}

\NewDocumentCommand \mc{ m }{%
    \mathcal{#1}%
}

\NewDocumentCommand \matD{ m m }{%
    \frac{{\rm D}#1}{{\rm D}#2}%
}

\NewDocumentCommand \D{ m m }{%
    \frac{{\rm d}#1}{{\rm d}#2}%
}

\NewDocumentCommand \kr{ m }{%
    \delta_{#1}
}

\NewDocumentCommand \tr{ m }{%
    \mbox{tr}_{\scs\omega} \left( {#1} \right)
}

\NewDocumentCommand \subm{ m }{%
    {#1}_{(m)}
}

\NewDocumentCommand \p{ m m }{%
    \frac{\partial #1}{\partial #2}%
}

\NewDocumentCommand \eref{ m }{%
    (\ref{eqn:#1})%
}

\NewDocumentCommand \sref{ m }{%
    \S\ref{sec:#1}%
}

\NewDocumentCommand \fref{ m }{%
    fig.~\ref{fig:#1}%
}

\NewDocumentCommand \Fref{ m }{%
    Figure~\ref{fig:#1}%
}

\NewDocumentCommand \an{ m }{%
    \langle {#1} \rangle%
}

\NewDocumentCommand \mbf{ m }{%
    \mathbf{#1}
}

\NewDocumentCommand \bbar{ m }{%
    \bar{\mbf{#1}}
}

\NewDocumentCommand \bs{ m }{%
    \boldsymbol{#1}
}

\NewDocumentCommand \tbbar{ m }{%
    \tilde{\bar{\mbf{#1}}}
}

\NewDocumentCommand \tbar{ m }{%
    \tilde{\bar{#1}}
}

\NewDocumentCommand \tmbf{ m }{%
    \tilde{\mbf{#1}}
}


\NewDocumentCommand \stateOne{ m }{
    \underline{\mbf{#1}}%
}

\NewDocumentCommand \stateTwo{ m o }{%
    \IfNoValueTF{#2}{\stateOne{#1}}
    {
    \stateOne{#1}\an{#2}%
    }
}

\NewDocumentCommand \s{ m o o }{%
    \IfNoValueTF{#3}{\stateTwo{#1}[#2]}%
    {
    \underline{\mbf{#1}}({#3})\an{#2}%
    }
}

\NewDocumentCommand \stateOneI{ m m }{
    \underline{#1}_{#2}%
}

\NewDocumentCommand \stateTwoI{ m m o }{%
    \IfNoValueTF{#3}{\stateOneI{#1}{#2}}
    {
    \stateOneI{#1}{#2}\an{#3}%
    }
}

% \NewDocumentCommand \si{ m m o o }{%
%     \IfNoValueTF{#4}{\stateTwoI{#1}{#2}[#3]}%
%     {
%         \stateOneI{#1}{#2}(\mbf{#4})\an{#3}
%     }
% }

\NewDocumentCommand \scs{ m }{
    \underline{#1}%
}

\NewDocumentCommand \ds{ m }{
    {\rm d}\scs{#1}%
}

\NewDocumentCommand \smag{ m }{
    \vert\s{#1}\vert%
}



\title{A finite-deformation formulation for partially saturated reacting porous media with exchange-potential-driven transport}

\author[1,2,3]{John T. Foster\footnote{Electronic address: \href{mailto:john.foster@utexas.edu}{john.foster@utexas.edu}}}
\author[1]{Syed Talha Tirmizi}
\author[3,4]{Marc Hesse}

\affil[1]{The Hildebrand Department of Petroleum \& Geosystems Engineering, The University of Texas at Austin}
\affil[2]{Department of Aerospace Engineering \& Engineering Mechanics, The University of Texas at Austin}
\affil[3]{The Oden Institute for Computational Engineering \& Science, The University of Texas at Austin}
\affil[4]{The Jackson School of Geosciences, The University of Texas at Austin}

\date{\today}

\begin{document}\sloppy
\maketitle


\section{Momentum balance}
\label{sec:momentum}

\subsection{Kinematics and phase structure}
The volume fractions satisfy
\begin{equation}
\phi_A + \phi_B + \phi_f + \phi_g = 1,
\label{eq:saturation}
\end{equation}
with total porosity $\phi=\phi_f+\phi_g$ and fluid saturation $S_f=\phi_f/\phi$. Each component map has $\mbf{F}_\xi=\partial\bs\chi_\xi/\partial\mbf{X}_\xi$ and $J_\xi=\det(\mbf{F}_\xi)$; solids share trajectory $\mbf{v}_A=\mbf{v}_B=\mbf{v}_s$ and $\mbf{F}_A=\mbf{F}_B=\mbf{F}_s$.

With multiplicative decomposition,
\begin{equation}
\mbf{F}_\xi=\mbf{A}_\xi\bar{\mbf{F}}_\xi, \qquad J_\xi=\alpha_\xi\bar J_\xi,
\end{equation}
and for incompressible solid phases,
\begin{equation}
\alpha_s=J_s,\qquad \bar{\mbf{F}}_s=J_s^{-1/3}\mbf{F}_s,\qquad \bar J_s\equiv1.
\label{eq:alpha_Js}
\end{equation}

\subsection{Pressure identifications from the variational principle}
From Appendix~\ref{app:variational}, thermodynamic pressures are
\begin{equation}
\bar p_\xi\equiv\bar\rho_\xi^2\frac{\partial\psi_\xi}{\partial\bar\rho_\xi},
\label{eq:pbar}
\end{equation}
with $\bar p_A=\bar p_B=-\lambda$. Fluid and gas pressures satisfy
\begin{align}
\bar p_f &= -\lambda-P_\phi^\star-p_c(1-S_f),\\
\bar p_g &= -\lambda-P_\phi^\star+p_cS_f,
\end{align}
where
\begin{equation}
p_c\equiv-\frac{1}{\phi}\frac{\partial\gamma_{fg}}{\partial S_f},
\qquad
P_\phi^\star\equiv\sum_{\xi=A,B}\rho_\xi\left.\frac{\partial\psi_\xi}{\partial\phi}\right|_{\mbf{F}_s}.
\end{equation}

\subsection{General gas-phase formulation with compressibility}
For a compressible gas phase without assuming atmospheric connectivity, the gas and fluid pressures remain related through capillarity:
\begin{align}
\bar p_f &= -\lambda-P_\phi^\star-p_c(1-S_f),\\
\bar p_g &= -\lambda-P_\phi^\star+p_cS_f.
\label{eq:pf_pg_general}
\end{align}
Hence the saturation-weighted phase pressure is
\begin{equation}
\bar p \equiv S_f\bar p_f + (1-S_f)\bar p_g = -\lambda - P_\phi^\star,
\label{eq:pbar_mix}
\end{equation}
so that
\begin{equation}
\lambda = -\bar p - P_\phi^\star.
\label{eq:lambda_from_pbar}
\end{equation}
The capillary pressure $p_c(S_f)$ remains bounded by capillary physics, ensuring physical closure of the saturation-pressure relationship.

\subsection{Component and mixture momentum}
\begin{align}
\rho_s\mbf{a}_s &= \nabla\cdot\left[(\bs\sigma_s')^\intercal+\lambda\mbf{I}+\phi P_\phi^\star\mbf{I}\right] + \rho_s\mbf{g}+\mbf{b}_s + c_f^+(\mbf{v}_s-\nabla\tau) + c_g^+(\mbf{v}_s-\nabla\tau),
\label{eq:mom_solid_component}\\
\rho_f\mbf{a}_f &= \rho_f\mbf{g}+\mbf{b}_f-\phi_f\nabla p_f + c_f^+(\nabla\tau-\mbf{v}_f),
\label{eq:mom_fluid_component}\\
\rho_g\mbf{a}_g &= \rho_g\mbf{g}+\mbf{b}_g-\phi_g\nabla p_g + c_g^+(\nabla\tau-\mbf{v}_g).
\label{eq:mom_gas_component}
\end{align}
Summing Eqs.~\eqref{eq:mom_solid_component}--\eqref{eq:mom_gas_component}, taking $\mbf{b}_s+\mbf{b}_f+\mbf{b}_g=\mbf{0}$, and using $\bar p = S_f\bar p_f+(1-S_f)\bar p_g = -\lambda-P_\phi^\star$ gives
\begin{equation}
\rho\mbf{a}=\nabla\cdot\left[(\bs\sigma_s')^\intercal-\bar p\mbf{I}-(1-\phi)P_\phi^\star\mbf{I}\right]+\rho\mbf{g}-c_f^+(\mbf{v}_f-\mbf{v}_s)-c_g^+(\mbf{v}_g-\mbf{v}_s).
\label{eq:mom_overall_main}
\end{equation}
Equivalently, since $\lambda=-\bar p-P_\phi^\star$,
\begin{equation}
\rho\mbf{a}=\nabla\cdot\left[(\bs\sigma_s')^\intercal+\lambda\mbf{I}+\phi P_\phi^\star\mbf{I}\right]+\rho\mbf{g}-c_f^+(\mbf{v}_f-\mbf{v}_s)-c_g^+(\mbf{v}_g-\mbf{v}_s).
\label{eq:mom_overall_lambda}
\end{equation}

\section{Constitutive modeling}
\label{sec:constitutive}

\subsection{Free-energy structure and distention energetics}
\begin{equation}
\psi_\xi(\mbf{F}_s,\phi)=\psi_\xi^{\mathrm{iso}}(\bar{\mbf{F}}_s)+\psi_\xi^{\mathrm{dis}}(J_s,\phi),\quad \xi\in\{A,B\},
\label{eq:psi_decomp}
\end{equation}
with $\psi_f(\bar\rho_f)$ and $\psi_g(\bar\rho_g)$.
Define
\begin{equation}
\bar P_\xi^\star=\left.\frac{\partial W_\xi^{\mathrm{dis}}}{\partial J_s}\right|_\phi+\frac{1-\phi}{J_s}\left.\frac{\partial W_\xi^{\mathrm{dis}}}{\partial\phi}\right|_{J_s}.
\end{equation}
Then admissible effective PK1 stress is
\begin{equation}
\mbf{P}_s'=J_s^{2/3}\sum_{\xi=A,B}\phi_\xi\left[\bar{\mbf{P}}_\xi'-\frac13(\bar{\mbf{P}}_\xi':\mbf{F}_s)\mbf{F}_s^{-T}\right]+J_s\sum_{\xi=A,B}\phi_\xi\bar P_\xi^\star\mbf{F}_s^{-T}.
\label{eq:PK1_total}
\end{equation}

\subsection{Specific model: neo-Hookean + logarithmic distention energy}
\begin{equation}
W_\xi^{\mathrm{iso}}=\frac{G_\xi}{2}(\bar{\mbf{F}}_s:\bar{\mbf{F}}_s-3),
\qquad
W_\xi^{\mathrm{dis}}=\frac{K_\xi^\star}{2(1-\phi)}(\ln J_s)^2.
\end{equation}
Hence
\begin{equation}
\mbf{P}_s' = \mu_{\mathrm{eff}}J_s^{1/3}\left[\mbf{F}_s-\frac13I_1\mbf{F}_s^{-T}\right] + K_{\mathrm{eff}}^\star\frac{\ln J_s}{J_s}\left(1+\frac12\ln J_s\right)\mbf{F}_s^{-T}.
\label{eq:PK1_specific}
\end{equation}

\subsection{Interfacial energy, capillary pressure, and saturation closure}

The capillary pressure is defined as the pressure difference between gas and fluid phases:
\begin{equation}
p_c \equiv p_g - p_f > 0.
\label{eq:pc_def_main}
\end{equation}
Interfacial energy determines this pressure variationally.  Writing the current
interfacial energy density as $\phi\gamma_{fg}(S_f)$, the primitive
volume-fraction derivatives are
\begin{align}
  \gamma_g &= \gamma_{fg}+S_f p_c, \notag\\
  \gamma_f &= \gamma_{fg}-(1-S_f)p_c,
\end{align}
with
\begin{equation}
  p_c = -\frac{\partial\gamma_{fg}}{\partial S_f}
      = \gamma_g-\gamma_f
      = \bar p_g-\bar p_f .
\end{equation}
The variational derivation in Appendix~\ref{sec:variational} therefore gives
the capillary-pressure condition directly from the water and gas
volume-fraction equations, without introducing an independent saturation
multiplier. We adopt a Brooks--Corey law written in terms of the effective liquid saturation
\begin{equation}
S_e \equiv \frac{S_f-S_r}{1-S_r-S_{nr}},
\qquad 0 < S_e \le 1,
\end{equation}
where $S_r$ and $S_{nr}$ are the residual wetting and nonwetting saturations. The capillary pressure is then
\begin{equation}
p_c(S_f) = p_d S_e^{-1/d} = p_d\left(\frac{S_f-S_r}{1-S_r-S_{nr}}\right)^{-1/d},
\label{eq:brooks_corey}
\end{equation}
where $p_d$ is the air-entry pressure and $d$ is the pore-size distribution index. Inverting this relationship gives
\begin{equation}
S_f(p_c) = S_r + (1-S_r-S_{nr})\left(\frac{p_d}{p_c}\right)^d,
\label{eq:Sf_of_pc}
\end{equation}
which is well-defined for $p_c \ge p_d > 0$. The liquid volume fraction then follows:
\begin{equation}
\phi_f = \phi \cdot S_f(p_c) = \phi\left[S_r + (1-S_r-S_{nr})\left(\frac{p_d}{p_c}\right)^d\right],
\label{eq:phi_f_of_pc}
\end{equation}
and the gas porosity is
\begin{equation}
\phi_g = \phi - \phi_f = \phi\left[(1-S_r) - (1-S_r-S_{nr})\left(\frac{p_d}{p_c}\right)^d\right].
\label{eq:phi_g_of_pc}
\end{equation}

Time derivatives propagate through the chain rule. From $p_c = p_g - p_f$:
\begin{equation}
\dot p_c = \dot p_g - \dot p_f,
\label{eq:pc_rate}
\end{equation}
and 
\begin{equation}
\dot S_f = \frac{\partial S_f}{\partial p_c}\dot p_c = -\frac{(1-S_r-S_{nr})d\,p_d^d}{p_c^{d+1}}\dot p_c.
\label{eq:Sf_rate}
\end{equation}
The time derivative of $\phi_f$ is:
\begin{equation}
\dot\phi_f = \dot\phi \cdot S_f + \phi \cdot \dot S_f.
\label{eq:phi_f_rate}
\end{equation}
These chain-rule dependencies are essential for consistency with the solved variables $(\dot p_f, \dot p_g, \dot\phi)$.

\subsection{Relative flux, reaction kinetics, and mass-source identification}

The relative fluid flux (reference frame) is:
\begin{equation}
\mbf{w}=\rho_f\tilde{\mbf{k}}\left[\rho_f\mbf{g}-\phi_f\nabla \bar p_f + c_f^+(\nabla\tau-\mbf{v}_s)\right],
\label{eq:darcy_flux}
\end{equation}
with reference-frame flux $\mbf{W}=J_s\mbf{F}_s^{-1}\mbf{w}$. The term $c_f^+(\nabla\tau-\mbf{v}_s)$ arises from the exchange potential and the mass-source constraint, not as an independent constitutive assumption.

The reaction kinetics are taken in the saturation-weighted conversion form
\begin{equation}
\left.\frac{\partial c_\xi}{\partial t}\right|_{\mbf{X}_s}=\nu_\xi R, \qquad R=m_B r a_0 \, S_f(1-\xi),
\label{eq:reaction_rate}
\end{equation}
where
\begin{equation}
  \xi \equiv \frac{\phi_B}{1-\phi},
  \qquad
  1-\xi = \frac{\phi_A}{1-\phi},
\end{equation}
is the local solid conversion variable, and $\nu_\xi$ are stoichiometric coefficients satisfying the constituent mass balance. In this formulation the implemented solid source terms are obtained directly from the stoichiometric rate per unit reference volume as
\begin{equation}
  c_A^+ = \frac{\nu_A R}{J_s}, \qquad c_B^+ = \frac{\nu_B R}{J_s},
\label{eq:c_AB_from_kinetics}
\end{equation}
with $\nu_A<0$ and $\nu_B>0$ for reactant consumption and product generation, respectively.

For the gas-off closure used in the present source-term implementation, the fluid source is taken as
\begin{equation}
  c_f^+ = -(c_A^+ + c_B^+),
\end{equation}
which is the no-gas limit of overall constituent mass conservation. In the fully trapped-gas formulation one instead has the more general combined constraint
\begin{equation}
  c_f^+ + c_g^+ = -\dot\phi + (1-\phi)\mathrm{tr}(\mbf{L}_s) - \frac{c_A^+}{\bar\rho_A} - \frac{c_B^+}{\bar\rho_B},
\end{equation}
but the current code path uses the gas-off closure above for the reaction-source partition.

\section{Mass balances}
\label{sec:mass_balances}

\subsection{Four-component mass conservation}
\begin{subequations}
\begin{align}
\frac{\partial(\phi_f\bar\rho_f)}{\partial t}+\nabla\cdot(\phi_f\bar\rho_f\mbf{v}_f)&=c_f^+,\label{eq:mb_f}\\
\frac{\partial(\phi_g\bar\rho_g)}{\partial t}+\nabla\cdot(\phi_g\bar\rho_g\mbf{v}_g)&=c_g^+,\label{eq:mb_g}\\
\frac{\partial\phi_A}{\partial t}+\nabla\cdot(\phi_A\mbf{v}_s)&=\frac{c_A^+}{\bar\rho_A},\label{eq:mb_A}\\
\frac{\partial\phi_B}{\partial t}+\nabla\cdot(\phi_B\mbf{v}_s)&=\frac{c_B^+}{\bar\rho_B}.\label{eq:mb_B}
\end{align}
\end{subequations}

\subsection{Constraint on mass-production rates}
The volume fraction constraint $\phi_A + \phi_B + \phi_f + \phi_g = 1$ and geometric incompressibility of solids ($\alpha_s = J_s$) together imply:
\begin{equation}
\sum_{\xi=A,B,f,g} c_\xi^+ = 0.
\label{eq:mass_balance_constraint}
\end{equation}
Thus $c_\xi^+$ is not independently specified but is derived from the volume-fraction and solid-mass constraints. Specifically, from the solid porosity kinematic equation and the volume-fraction constraint:
\begin{equation}
c_f^+ + c_g^+ = -\dot\phi + (1-\phi)\mathrm{tr}(\mbf{L}_s) - \frac{c_A^+}{\bar\rho_A} - \frac{c_B^+}{\bar\rho_B}.
\label{eq:c_fg_from_constraint}
\end{equation}
Since the reaction rates $c_A^+$ and $c_B^+$ are driven by kinetics (external), and we solve for $\dot\phi$ from the solid-mass balance, the combined flux $c_f^+ + c_g^+$ is geometrically determined. The individual partition between $c_f^+$ and $c_g^+$ follows from the saturation kinetics and capillary coupling.

\subsection{Solid porosity balance}
With $\phi_B=1-\phi-\phi_A$, integrating the volume-fraction and solid-mass constraints gives:
\begin{equation}
-\dot\phi + (1-\phi)\,\mathrm{tr}(\mbf{L}_s) - \frac{c_A^+}{\bar\rho_A}-\frac{c_B^+}{\bar\rho_B}=0.
\label{eq:solid_mass_phi}
\end{equation}
This equation is solved for $\dot\phi$ (or equivalently, the rate of porosity change) given the reaction rates and skeleton deformation.

\subsection{Fluid and gas equations of state}

\subsubsection{Fluid phase (liquid water)}
For the liquid phase, assume linear compressibility:
\begin{equation}
\dot{\bar\rho}_f=\frac{\bar\rho_f}{K_f}\dot{\bar p}_f,
\label{eq:fluid_eos}
\end{equation}
where $K_f$ is the fluid bulk modulus. This is a constituent closure and is solved as one of the nine weak forms.

\subsubsection{Gas phase (ideal gas with trapped-gas closure)}
For the gas phase, assume ideal-gas behavior under isothermal conditions:
\begin{equation}
\bar\rho_g = \frac{M_g}{RT}\bar p_g,
\label{eq:ideal_gas_eos_intrinsic}
\end{equation}
where $M_g$ is the molar mass, $R$ is the universal gas constant, and $T$ is held constant. In the trapped-gas saturation formulation we do not solve a separate gas mass balance. Instead, we impose
\begin{equation}
  c_g^+ = 0,
  \qquad
  \mbf Q = \mbf 0,
  \label{eq:trapped_gas_assumption}
\end{equation}
and recover the gas pressure and density diagnostically from capillarity,
\begin{equation}
  \bar p_g = \bar p_f + p_c(S_f),
  \qquad
  \bar\rho_g = \frac{M_g}{RT}\left[\bar p_f + p_c(S_f)\right].
  \label{eq:pg_rhog_from_pf_S}
\end{equation}
Thus the gas phase acts as a trapped compressible storage phase with no independent reactive source term.

\subsubsection{Saturation equation from the fluid mass balance}
We retain the same notation used in the ReactingMixture paper and write
\begin{equation}
  S_f = \frac{\phi_f}{\phi},
  \qquad
  \phi_f = \phi S_f,
  \qquad
  \phi_g = \phi(1-S_f).
  \label{eq:Sf_definition_mass}
\end{equation}
Starting from the pulled-back fluid mass balance,
\begin{equation}
  J_s\frac{\partial(\phi_f\bar\rho_f)}{\partial t} + \phi_f\bar\rho_f\frac{\partial J_s}{\partial t} + \nabla_{\mbf{X}_s}\cdot\mbf W = J_s c_f^+,
  \label{eq:sat_fluid_start}
\end{equation}
and substituting $\phi_f = \phi S_f$ gives
\begin{equation}
  J_s\frac{\partial(\phi S_f \bar\rho_f)}{\partial t} + \phi S_f \bar\rho_f\frac{\partial J_s}{\partial t} + \nabla_{\mbf{X}_s}\cdot\mbf W = J_s c_f^+.
  \label{eq:sat_fluid_sub_1}
\end{equation}
Expanding the time derivative, dividing by $J_s$, and using
\begin{equation}
  \frac{\partial J_s}{\partial t} = J_s\,\mbox{tr}(\dot{\mbf F}_s\mbf F_s^{-1}),
  \label{eq:Js_rate_sat}
\end{equation}
gives
\begin{equation}
  \bar\rho_f\phi\dot S_f + \bar\rho_f S_f\dot\phi + \phi S_f\dot{\bar\rho}_f + \phi S_f\bar\rho_f\,\mbox{tr}(\dot{\mbf F}_s\mbf F_s^{-1}) + \frac{1}{J_s}\nabla_{\mbf{X}_s}\cdot\mbf W = c_f^+.
  \label{eq:sat_fluid_sub_3}
\end{equation}
Now use the solid mass balance,
\begin{equation}
  \dot\phi = (1-\phi)\,\mbox{tr}(\dot{\mbf F}_s\mbf F_s^{-1}) - \frac{c_A^+}{\bar\rho_A} - \frac{c_B^+}{\bar\rho_B},
  \label{eq:sat_solid_sub_2}
\end{equation}
and the fluid equation of state,
\begin{equation}
  \dot{\bar\rho}_f = \frac{\bar\rho_f}{K_f}\dot{\bar p}_f,
  \label{eq:sat_fluid_eos_sub}
\end{equation}
to obtain the saturation equation
\begin{equation}
  \boxed{\phi\dot S_f + \frac{\phi_f}{K_f}\dot{\bar p}_f + S_f\,\mbox{tr}(\dot{\mbf F}_s\mbf F_s^{-1}) - S_f\left(\frac{c_A^+}{\bar\rho_A} + \frac{c_B^+}{\bar\rho_B}\right) + \frac{1}{J_s\bar\rho_f}\nabla_{\mbf{X}_s}\cdot\mbf W - \frac{c_f^+}{\bar\rho_f} = 0.}
  \label{eq:sat_from_fluid_eos}
\end{equation}
This is the equation solved for $S_f$.

\subsubsection{Overall pressure equation from fluid plus gas storage}
The pressure equation is obtained by adding the fluid and gas mass balances under the trapped-gas assumptions~\eqref{eq:trapped_gas_assumption}. Since $\mbf Q=\mbf 0$ and $c_g^+=0$, we have
\begin{equation}
  J_s\frac{\partial(\phi_f\bar\rho_f + \phi_g\bar\rho_g)}{\partial t} + \left(\phi_f\bar\rho_f + \phi_g\bar\rho_g\right)\frac{\partial J_s}{\partial t} + \nabla_{\mbf{X}_s}\cdot\mbf W = J_s c_f^+.
  \label{eq:overall_start}
\end{equation}
Define
\begin{equation}
  \hat\rho \equiv S_f\bar\rho_f + (1-S_f)\bar\rho_g,
  \label{eq:rhohat_def}
\end{equation}
so that $\phi_f\bar\rho_f + \phi_g\bar\rho_g = \phi\hat\rho$. Then Eq.~\eqref{eq:overall_start} becomes
\begin{equation}
  \hat\rho\,\dot\phi + \phi\dot{\hat\rho} + \phi\hat\rho\,\mbox{tr}(\dot{\mbf F}_s\mbf F_s^{-1}) + \frac{1}{J_s}\nabla_{\mbf{X}_s}\cdot\mbf W - c_f^+ = 0.
  \label{eq:overall_phi_rhohat}
\end{equation}
Furthermore,
\begin{equation}
  \dot{\hat\rho} = (\bar\rho_f-\bar\rho_g)\dot S_f + S_f\dot{\bar\rho}_f + (1-S_f)\dot{\bar\rho}_g,
  \label{eq:rhohat_dot_1}
\end{equation}
and from Eqs.~\eqref{eq:sat_fluid_eos_sub} and~\eqref{eq:pg_rhog_from_pf_S},
\begin{equation}
  \dot{\bar\rho}_g = \frac{M_g}{RT}\left(\dot{\bar p}_f + p_c'(S_f)\dot S_f\right).
  \label{eq:rhog_dot_pfSf}
\end{equation}
Therefore,
\begin{equation}
  \dot{\hat\rho} = \left[S_f\frac{\bar\rho_f}{K_f} + (1-S_f)\frac{M_g}{RT}\right]\dot{\bar p}_f + \left[\bar\rho_f-\bar\rho_g + (1-S_f)\frac{M_g}{RT}p_c'(S_f)\right]\dot S_f.
  \label{eq:rhohat_dot_2}
\end{equation}
Substituting Eqs.~\eqref{eq:sat_solid_sub_2} and~\eqref{eq:rhohat_dot_2} into Eq.~\eqref{eq:overall_phi_rhohat} yields the exact overall pressure equation,
\begin{equation}
  \boxed{C_p\dot{\bar p}_f + C_S\dot S_f + \hat\rho\,\mbox{tr}(\dot{\mbf F}_s\mbf F_s^{-1}) - \hat\rho\left(\frac{c_A^+}{\bar\rho_A} + \frac{c_B^+}{\bar\rho_B}\right) + \frac{1}{J_s}\nabla_{\mbf{X}_s}\cdot\mbf W - c_f^+ = 0,}
  \label{eq:overall_pressure_equation}
\end{equation}
where
\begin{align}
  C_p &\equiv \phi\left[S_f\frac{\bar\rho_f}{K_f} + (1-S_f)\frac{M_g}{RT}\right],
  \label{eq:Cp_def}\\
  C_S &\equiv \phi\left[\bar\rho_f-\bar\rho_g + (1-S_f)\frac{M_g}{RT}p_c'(S_f)\right].
  \label{eq:Cs_def}
\end{align}
Equation~\eqref{eq:overall_pressure_equation} is solved for $\bar p_f$, while Eq.~\eqref{eq:sat_from_fluid_eos} is solved for $S_f$.

\section{System of eight unknowns and closures}
\label{sec:system_definition}

The trapped-gas saturation formulation solves for eight unknowns in the reference configuration:
\begin{enumerate}
  \item Solid displacements: $u_x, u_y$ (2 DOFs)
  \item Liquid pressure: $\bar p_f$ (1 DOF)
  \item Saturation: $S_f$ (1 DOF)
  \item Total porosity: $\phi$ (1 DOF)
  \item Reactant fraction: $\phi_A$ (1 DOF)
  \item Exchange potential: $\tau$ (1 DOF)
  \item Intrinsic liquid density: $\bar\rho_f$ (1 DOF)
\end{enumerate}
The gas pressure and gas density are recovered diagnostically from Eq.~\eqref{eq:pg_rhog_from_pf_S}.

All remaining variables are determined by closures:
\begin{itemize}
  \item Deformation gradient: $\mbf{F}_s = \mbf{I} + \nabla u$ (kinematic)
  \item Jacobian: $J_s = \det(\mbf{F}_s)$ (kinematic)
  \item Product fraction: $\phi_B = (1-\phi) - \phi_A$ (volume-fraction constraint)
  \item Liquid porosity: $\phi_f = \phi S_f$ (definitional)
  \item Gas porosity: $\phi_g = \phi(1-S_f)$ (definitional)
  \item Capillary pressure: $p_c = p_c(S_f)$ (capillary closure)
  \item Gas pressure: $\bar p_g = \bar p_f + p_c(S_f)$ (diagnostic)
  \item Gas density: $\bar\rho_g = \frac{M_g}{RT}[\bar p_f + p_c(S_f)]$ (diagnostic)
  \item Mass sources: $c_A^+, c_B^+, c_f^+$ with $c_g^+=0$
  \item Bulk densities: $\rho_\xi = \phi_\xi \bar\rho_\xi$ (definitional)
  \item Mixture density: $\rho = \sum \rho_\xi$ (definitional)
  \item Velocity gradients: $\mbf{L}_s = \dot{\mbf{F}}_s \mbf{F}_s^{-1}$ (kinematic)
\end{itemize}

\section{Exchange potential evolution}
\label{sec:tau}

\begin{equation}
\left.\frac{\partial\tau}{\partial t}\right|_{\mbf{X}_s}=\frac12\mbf{v}_s\cdot\mbf{v}_s+\phi_A\left.\frac{\partial\psi_A^{\mathrm{dis}}}{\partial\phi}\right|_{J_s}-\psi_A+\frac{\lambda}{\bar\rho_A}.
\label{eq:tau_evolution}
\end{equation}
When $P_\phi^\star\to0$, this recovers the prior three-phase form.

\section{The role of the exchange potential}
\label{sec:tau_role}

\subsection{Compressible gas and bounded suction}
With compressible gas and capillary closure, the gas pressure $\bar p_g$ evolves according to both mass balance and ideal-gas equation, while fluid pressure $\bar p_f$ remains bounded by capillary physics: $|\nabla \bar p_f| \le p_c(S_f)|\nabla S_f|$.

\subsection{Exchange-potential gradient as thermodynamic forcing}
Define
\begin{equation}
\varepsilon\equiv\frac{|c_f^+\nabla\tau|}{\phi_f|\nabla \bar p_f|}.
\label{eq:epsilon}
\end{equation}
Near active fronts, free-energy contrasts can make $\varepsilon>1$, which is especially important in regions of limited capillary-pressure gradients.

\section{Weak forms: the trapped-gas saturation formulation}
\label{sec:weak_forms}

\subsection{System summary}
The following weak forms determine the eight unknowns $(u_x, u_y, \bar p_f, S_f, \phi, \phi_A, \tau, \bar\rho_f)$, with the momentum balance contributing two component equations. The trapped-gas pressure and density are recovered diagnostically.

\paragraph{(1) Momentum balance}
\begin{equation}
\int_{\Omega_0}\left[\mbf{P}_s':\nabla_{\mbf{X}}\delta\mbf{u}-\bar p\,J_s\mbf{F}_s^{-T}:\nabla_{\mbf{X}}\delta\mbf{u} -(1-\phi)P_\phi^\star J_s\mbf{F}_s^{-T}:\nabla_{\mbf{X}}\delta\mbf{u} -J_s\rho\mbf{g}\cdot \delta\mbf{u} +\frac{c_f^+}{\phi_f\bar\rho_f}\mbf{F}_s\mbf{W}\cdot \delta\mbf{u}\right]dV_0=0.
\label{eq:wf_mom}
\end{equation}
Here we used the current-frame relative mass flux $\mbf w=\phi_f\bar\rho_f(\mbf v_f-\mbf v_s)$ together with the pull-back relation $\mbf W=J_s\mbf F_s^{-1}\mbf w$, so that
\begin{equation}
  \mbf v_f-\mbf v_s = \frac{\mbf F_s\mbf W}{J_s\phi_f\bar\rho_f}.
\end{equation}
Therefore the reaction-momentum contribution $-J_s c_f^+(\mbf v_f-\mbf v_s)\cdot\delta\mbf u$ becomes $-(c_f^+/(\phi_f\bar\rho_f))\mbf F_s\mbf W\cdot\delta\mbf u$ inside the pulled-back weak form, and after moving all terms to the left this yields the positive contribution shown above.

\paragraph{(2) Overall pressure equation}
\begin{equation}
\int_{\Omega_0}\left[\delta\bar p_f\,J_s\left(C_p\dot{\bar p}_f + C_S\dot S_f + \hat\rho\,\mathrm{tr}(\mbf{L}_s) - \hat\rho\left(\frac{c_A^+}{\bar\rho_A} + \frac{c_B^+}{\bar\rho_B}\right) - c_f^+\right) - \nabla_{\mbf{X}}\delta\bar p_f\cdot\mbf W\right]dV_0=0.
\label{eq:wf_pressure}
\end{equation}

\paragraph{(3) Saturation equation}
\begin{equation}
\int_{\Omega_0}\left[\delta S_f\,J_s\left(\phi\dot S_f + \frac{\phi_f}{K_f}\dot{\bar p}_f + S_f\,\mathrm{tr}(\mbf{L}_s) - S_f\left(\frac{c_A^+}{\bar\rho_A} + \frac{c_B^+}{\bar\rho_B}\right) - \frac{c_f^+}{\bar\rho_f}\right) - \nabla_{\mbf{X}}\delta S_f\cdot\frac{\mbf W}{\bar\rho_f}\right]dV_0=0.
\label{eq:wf_saturation}
\end{equation}

\paragraph{(4) Solid porosity balance}
\begin{equation}
\int_{\Omega_0}\left[-\delta\phi\,J_s\dot\phi + \delta\phi\,J_s(1-\phi)\mathrm{tr}(\mbf{L}_s)-\delta\phi\,J_s\frac{c_A^+}{\bar\rho_A}-\delta\phi\,J_s\frac{c_B^+}{\bar\rho_B}\right]dV_0=0.
\label{eq:wf_solid}
\end{equation}

\paragraph{(5) Reactant mass balance}
\begin{equation}
\int_{\Omega_0}\left[\delta\phi_A\,J_s\dot\phi_A+\delta\phi_A\,J_s\phi_A\mathrm{tr}(\mbf{L}_s)-\delta\phi_A\,J_s\frac{c_A^+}{\bar\rho_A}\right]dV_0=0.
\label{eq:wf_reactant}
\end{equation}

\paragraph{(6) Exchange potential evolution}
\begin{equation}
\int_{\Omega_0}\delta\tau\,J_s\left[\dot\tau-\frac12\mbf{v}_s\cdot\mbf{v}_s-\phi_A\left.\frac{\partial\psi_A^{\mathrm{dis}}}{\partial\phi}\right|_{J_s}+\psi_A-\frac{\lambda}{\bar\rho_A}\right]dV_0=0.
\label{eq:wf_tau}
\end{equation}

\paragraph{(7) Fluid equation of state}
\begin{equation}
\int_{\Omega_0}\delta\bar\rho_f\,J_s\left[\dot{\bar\rho}_f-\frac{\bar\rho_f}{K_f}\dot{\bar p}_f\right]dV_0=0.
\label{eq:wf_eos_f}
\end{equation}

\paragraph{(8) Diagnostic gas closure}
After solving Eqs.~\eqref{eq:wf_pressure}--\eqref{eq:wf_eos_f}, the trapped-gas quantities follow pointwise from
\begin{equation}
  \bar p_g = \bar p_f + p_c(S_f), \qquad \bar\rho_g = \frac{M_g}{RT}\left[\bar p_f + p_c(S_f)\right], \qquad \phi_g = \phi(1-S_f).
\label{eq:diag_pc}
\end{equation}

\section{Conclusions}
\label{sec:conclusions}

We reorganized the four-phase framework to match the three-phase manuscript architecture while retaining partial-saturation physics, distention energy, and a full variational/thermodynamic closure.

\section*{Acknowledgments}
The authors thank Syed Talha Tirmizi and Emmanuel Ikpesu for implementation and verification support in MOOSE.

\appendix
\section{Variational Derivation}
\label{app:variational}
\label{sec:variational}

We derive the governing equations from Hamilton's extended principle, following Bedford \& Drumheller~\cite{drumheller1980thermomechanical} and Drumheller~\cite{drumheller2000}, adapted to a partially saturated four-phase system with porosity-dependent solid free energies.

Each component has volume fraction~$\phi_\xi$, intrinsic density~$\bar\rho_\xi$, bulk density $\rho_\xi = \phi_\xi\bar\rho_\xi$, velocity~$\mbf{v}_\xi$, deformation gradient~$\mbf{F}_\xi$, and Jacobian $J_\xi = \det(\mbf{F}_\xi)$.  The solid phases share a common trajectory ($\mbf{v}_A = \mbf{v}_B = \mbf{v}_s$, $\mbf{F}_A = \mbf{F}_B = \mbf{F}_s$, $J_A = J_B = J_s$).  The total porosity is $\phi = \phi_f + \phi_g$, the fluid saturation is $S_f = \phi_f/\phi$, and the mass added per unit reference volume to component~$\xi$ by reactions is~$c_\xi$, with rate per unit current volume $c^+_\xi = \dot{c}_\xi / J_\xi$.

The free energies of the solid phases are written as $\psi_A(\mbf{F}_s, \phi)$ and $\psi_B(\mbf{F}_s, \phi)$, retaining the full deformation gradient and the current porosity as independent arguments.  The multiplicative decomposition $\mbf{F}_s = \mbf{A}_s \bar{\mbf F}_s$ is \emph{not} imposed at this stage; it enters later in the Coleman--Noll procedure (Appendix~\ref{sec:coleman}) when a specific constitutive form is adopted.

\subsection{Variational principle and its terms}

The extended Hamilton's principle is
\begin{equation*}
  \int_{t_1}^{t_2}\left[\delta(T - U) + \delta W + \sum_k \delta C_k\right] dt = 0,
\end{equation*}
with kinetic energy, internal energy, virtual work, and three constraints:
\begin{align*}
  T &= \sum_\xi \int_{\mc{B}} \tfrac{1}{2}\rho_\xi\,\mbf{v}_\xi\cdot\mbf{v}_\xi \,\dv,  \\
  U &= \int_{\mc{B}_{s0}} \left[\rho_{A0}\,\psi_A(\mbf{F}_s, \phi) + \rho_{B0}\,\psi_B(\mbf{F}_s, \phi)\right] dV_s \\
    &\qquad + \int_{\mc{B}_{f0}} \rho_{f0}\,\psi_f(\bar\rho_f)\,dV_f + \int_{\mc{B}_{g0}} \rho_{g0}\,\psi_g(\bar\rho_g)\,dV_g + \int_{\mc{B}} \gamma_{fg}(S_f)\,\dv,  \\
  \delta W &= \sum_\xi \int_{\mc{B}} \left[(\rho_\xi\mbf{g} + \mbf{b}_\xi)\cdot\delta\mbf{x}_\xi + \mc{L}_\xi\,\frac{\tilde\delta c_\xi}{J_\xi}\right] \,\dv,  \\
  C_1 &= \int_{\mc{B}} \lambda\left(\sum_\xi\phi_\xi - 1\right) \,\dv,  \\
  C_2 &= \sum_\xi \int_{\mc{B}_{\xi 0}} \mu_\xi\left(J_\xi - \frac{\phi_{\xi 0}\bar\rho_{\xi 0} + c_\xi}{\phi_\xi\bar\rho_\xi}\right) dV_\xi,  \\
  C_3 &= \int_{\mc{B}} \tau\sum_\xi c^+_\xi \,\dv.
\end{align*}
The Helmholtz free energies are $\psi_A(\mbf{F}_s, \phi)$, $\psi_B(\mbf{F}_s, \phi)$, $\psi_f(\bar\rho_f)$, and $\psi_g(\bar\rho_g)$.  The interaction forces satisfy $\sum_\xi \mbf{b}_\xi = \mbf{0}$, and $\mc{L}_\xi$ is a generalized force conjugate to mass transfer (identified later via Coleman--Noll).  The Lagrange multipliers~$\lambda$, $\mu_\xi$, and~$\tau$ enforce the volume fraction constraint, component mass conservation, and mixture mass balance, respectively.

Most of the required variations are detailed in Bedford \& Drumheller~\cite{drumheller1980thermomechanical} and will be omitted here; however, there are some new terms in this work that are derived below.

\subsection{Variation of the internal energy}

The solid contribution is
\begin{align}
  \delta U_s &= \int_{\mc{B}_{s0}} \sum_{\xi=A,B} \rho_{\xi 0} \left[\left.\p{\psi_\xi}{F_{kl}}\right\vert_\phi \delta F_{kl} + \left.\p{\psi_\xi}{\phi}\right\vert_{\mbf{F}_s} \delta\phi \right] dV_s.
  \label{eq:deltaU_solid}
\end{align}
The variation $\delta F_{kl} = \partial(\delta x_{s,k})/\partial X_{s,l}$ produces the standard stress divergence upon integration by parts.  Transforming to the current configuration gives
\begin{equation}
  \int_{\mc{B}_{0}} \rho_{\xi 0}\,\left.\p{\psi_\xi}{F_{kl}}\right\vert_\phi \delta F_{kl}\, dV_s = -\int_{\mc{B}} \nabla_{\mbf{x}} \cdot (\bs\sigma'_\xi)^T \cdot \delta\mbf{x}_s \,\dv + \text{boundary terms},
  \label{eq:stress_variation}
\end{equation}
where the Terzaghi effective Cauchy stress, i.e.\ the stress obtained with the
mobile-phase intrinsic densities held fixed, is defined through
\begin{equation}
  \bs\sigma'_\xi = \frac{1}{J_s} \rho_{\xi 0}\,\p{\psi_\xi}{\mbf{F}_s}\bigg\vert_\phi\, \mbf{F}_s^T.
  \label{eq:var_cauchy_def}
\end{equation}
The notation $|_\phi$ emphasizes that this derivative is taken at fixed porosity.

\subsection{Porosity variation via the inter-component identity}

The porosity~$\phi = \phi_f + \phi_g$ is a spatial field naturally associated with the fluid and gas components, yet it appears as an argument of the solid free energies $\psi_\xi(\mbf{F}_s, \phi)$ in the integral over the solid reference configuration.  Since this integral is taken at fixed~$\mbf{X}_s$, we require the variation of~$\phi$ at fixed~$\mbf{X}_s$.  Using the inter-component variation identity (equation~(1.26) of Bedford \& Drumheller~\cite{bedford1983theories})
\begin{equation}
  \left.\delta\Gamma_{(\zeta)}\right|_{\mbf{X}_{(\xi)}} = \delta\Gamma_{(\zeta)} - \nabla_{\mbf{x}}\Gamma_{(\zeta)} \cdot \left(\delta\mbf{x}_{(\zeta)} - \delta\mbf{x}_{(\xi)}\right),
  \label{eq:BD_identity}
\end{equation}
we obtain
\begin{align}
  \left.\delta\phi\right|_{\mbf{X}_s}
  &= \delta\phi_f + \delta\phi_g - \nabla_{\mbf{x}}\phi_f \cdot \left(\delta\mbf{x}_f - \delta\mbf{x}_s\right) - \nabla_{\mbf{x}}\phi_g \cdot \left(\delta\mbf{x}_g - \delta\mbf{x}_s\right) \notag \\
  &= \delta\phi_f + \delta\phi_g - \nabla_{\mbf{x}}\phi_f \cdot \delta\mbf{x}_f - \nabla_{\mbf{x}}\phi_g \cdot \delta\mbf{x}_g + \nabla_{\mbf{x}}\phi \cdot \delta\mbf{x}_s.
  \label{eq:delta_phi}
\end{align}
Defining
\begin{equation}
  P^\star_\phi \equiv \sum_{\xi=A,B} \rho_\xi\,\left.\p{\psi_\xi}{\phi}\right\vert_{\mbf{F}_s},
  \label{eq:Pstar_phi}
\end{equation}
and substituting~\eqref{eq:delta_phi} into~\eqref{eq:deltaU_solid}, the porosity dependent part of the solid energy variation becomes:
\begin{align}
  \int_{\mc{B}_{0}} \sum_{\xi=A,B} \rho_{\xi 0}\,\left.\p{\psi_\xi}{\phi}\right\vert_{\mbf{F}_s}\, \left.\delta\phi\right|_{\mbf{X}_s}\, dV_s
  &= \int_{\mc{B}} P^\star_\phi\, \delta\phi_f \,\dv + \int_{\mc{B}} P^\star_\phi\, \delta\phi_g \,\dv \notag \\
  &\quad - \int_{\mc{B}} P^\star_\phi\, \nabla_{\mbf{x}}\phi_f \cdot \delta\mbf{x}_f \,\dv - \int_{\mc{B}} P^\star_\phi\, \nabla_{\mbf{x}}\phi_g \cdot \delta\mbf{x}_g \,\dv \notag \\
  &\quad + \int_{\mc{B}} P^\star_\phi\, \nabla_{\mbf{x}}\phi \cdot \delta\mbf{x}_s \,\dv.
  \label{eq:phi_variation_split}
\end{align}

\subsection{Variation of the interfacial energy}

With two mobile phases (liquid water and gas), let
\[
  \phi = \phi_f+\phi_g,\qquad
  S_f=\frac{\phi_f}{\phi},\qquad
  S_g=1-S_f=\frac{\phi_g}{\phi}.
\]
The saturations are used only as ratios of the primitive volume-fraction fields;
the independent variations remain $\delta\phi_f$ and $\delta\phi_g$.  We write
the current interfacial energy density as
\[
  U_{fg}=\int_{\mc B}\phi\,\gamma_{fg}(S_f)\,\dv .
\]
Define the liquid--gas capillary pressure by
\begin{equation}
  p_c \equiv -\frac{\partial\gamma_{fg}}{\partial S_f} > 0.
  \label{eq:pc_def}
\end{equation}
Using $\delta\phi=\delta\phi_f+\delta\phi_g$ and
\[
  \delta S_f=\frac{\delta\phi_f-S_f\delta\phi}{\phi},
\]
the algebraic part of the interfacial variation is
\begin{align}
  \delta\!\left(\phi\gamma_{fg}\right)
  &= \gamma_{fg}\,\delta\phi
     + \frac{\partial\gamma_{fg}}{\partial S_f}
       \left(\delta\phi_f-S_f\delta\phi\right) \notag\\
  &= \gamma_f\,\delta\phi_f+\gamma_g\,\delta\phi_g,
\end{align}
where
\begin{align}
  \gamma_g &= \gamma_{fg}+S_f p_c, \notag\\
  \gamma_f &= \gamma_g-p_c
            = \gamma_{fg}-(1-S_f)p_c .
  \label{eq:primitive_gamma_derivatives}
\end{align}
Consequently the capillary pressure is the difference of primitive interfacial
potentials,
\begin{equation}
  p_c \equiv \bar p_g-\bar p_f \equiv \gamma_g-\gamma_f .
  \label{eq:pc_gamma_difference}
\end{equation}

The spatial fields must also be varied with their transported parts.  For the
phase displacement $\delta\mbf{x}_\xi$,
\[
  \delta_x\phi_\xi=-\nabla_{\mbf x}\cdot(\phi_\xi\delta\mbf{x}_\xi),
\]
so, after integration by parts, the complete interfacial-energy contribution is
\begin{align}
  \delta U_{fg}
  &= \int_{\mc{B}} \left(\gamma_f\,\delta\phi_f
      + \gamma_g\,\delta\phi_g\right) \,\dv \notag \\
  &\quad + \int_{\mc{B}} \left(\phi_f\nabla_{\mbf{x}}\gamma_f\cdot\delta\mbf{x}_f
      + \phi_g\nabla_{\mbf{x}}\gamma_g\cdot\delta\mbf{x}_g\right) \,\dv
  + \text{boundary terms}.
  \label{eq:deltaU_fg}
\end{align}
This form retains the gradient terms generated by the spatial dependence of the
interfacial potentials; they are eliminated only after the volume-fraction
Euler--Lagrange equations are used.

\subsection{Collected variational principle}

Evaluating all variations and collecting by independent variations $\delta\mbf{x}_\xi$, $\delta\phi_\xi$, $\delta\bar\rho_\xi$, and $\delta c_\xi$ gives
\begin{align}
  0 = \sum_\xi \int_{\mc{B}} \bigg\{
  &\bigg[-\rho_\xi\mbf{a}_\xi + \nabla_{\mbf{x}}\cdot(\bs\sigma'_\xi)^T + \rho_\xi\mbf{g} + \mbf{b}_\xi
  - \nabla_{\mbf{x}}\left(\mu_\xi J_\xi\right) - \lambda\nabla_{\mbf{x}}\phi_\xi + c^+_\xi\nabla_{\mbf{x}}\tau
  \bigg]\cdot\delta\mbf{x}_\xi \notag \\
  &+ \bigg[\frac{\mu_\xi J_\xi}{\phi_\xi} + \lambda\bigg] \,\delta\phi_\xi
  + \bigg[\frac{\mu_\xi J_\xi}{\bar\rho_\xi} - \rho_\xi\p{\psi_\xi}{\bar\rho_\xi}\bigg] \,\delta\bar\rho_\xi \notag \\
  &+ \bigg[\frac{1}{2}\mbf{v}_\xi\cdot\mbf{v}_\xi + \mc{L}_\xi - \frac{D_\xi\tau}{Dt} - \frac{\mu_\xi J_\xi}{\rho_\xi}\bigg]\frac{\delta c_\xi}{J_\xi}
  \bigg\} \,\dv \notag \\
  &+ \int_{\mc{B}} P^\star_\phi\, \delta\phi_f \,\dv + \int_{\mc{B}} P^\star_\phi\, \delta\phi_g \,\dv
  + \int_{\mc{B}} P^\star_\phi\, \nabla_{\mbf{x}}\phi \cdot \delta\mbf{x}_s \,\dv \notag \\
  &- \int_{\mc{B}} P^\star_\phi\, \nabla_{\mbf{x}}\phi_f \cdot \delta\mbf{x}_f \,\dv - \int_{\mc{B}} P^\star_\phi\, \nabla_{\mbf{x}}\phi_g \cdot \delta\mbf{x}_g \,\dv \notag \\
  &- \int_{\mc{B}} \left(\gamma_f\,\delta\phi_f
      + \gamma_g\,\delta\phi_g
      + \phi_f\nabla_{\mbf{x}}\gamma_f\cdot\delta\mbf{x}_f
      + \phi_g\nabla_{\mbf{x}}\gamma_g\cdot\delta\mbf{x}_g\right) \,\dv.
  \label{eq:collected}
\end{align}
This is the fully collected first variation.  The final two lines contain the new contributions from the porosity-dependent solid energy and the fluid--gas interfacial energy.

\subsection{Euler--Lagrange equations}

Since the variations are independent and arbitrary, each coefficient in~\eqref{eq:collected} must vanish.

\paragraph{Pressure identifications.}  From the coefficient of $\delta\bar\rho_\xi$ we identify the thermodynamic phase pressures
\begin{equation}
  \bar{p}_\xi \equiv \bar\rho_\xi^2\frac{\partial\psi_\xi}{\partial\bar\rho_\xi}, \qquad \xi \in \{f,g\},
  \label{eq:pbar_def}
\end{equation}
so that
\begin{equation}
  \frac{\mu_f J_f}{\phi_f} = \bar p_f, \qquad \frac{\mu_g J_g}{\phi_g} = \bar p_g.
  \label{eq:mu}
\end{equation}
From $\delta\phi_A$ and $\delta\phi_B$,
\begin{equation}
  \frac{\mu_A J_s}{\phi_A} = -\lambda, \qquad \frac{\mu_B J_s}{\phi_B} = -\lambda.
\end{equation}
From $\delta\phi_f$ and $\delta\phi_g$, including the $P^\star_\phi$ and
interfacial-energy contributions,
\begin{align}
  \bar p_f &= -\lambda - P^\star_\phi + \gamma_f,
  \label{eq:pf}
  \\
  \bar p_g &= -\lambda - P^\star_\phi + \gamma_g.
  \label{eq:pg}
\end{align}
Subtracting~\eqref{eq:pf} from~\eqref{eq:pg} gives the capillary-pressure
condition
\begin{equation}
  \bar p_g-\bar p_f = \gamma_g-\gamma_f = p_c .
  \label{eq:pc_from_phi_variations}
\end{equation}
Therefore
\begin{equation}
  \bar p_E \equiv S_f\bar p_f + (1-S_f)\bar p_g - \gamma_{fg}
  = -\lambda - P^\star_\phi,
  \label{eq:lambda}
\end{equation}
where $S_f\gamma_f+S_g\gamma_g=\gamma_{fg}$ for the specialized density
$\phi\gamma_{fg}(S_f)$.  Thus, when the solid energy has no explicit porosity
dependence, $\bar p_E=-\lambda$, matching the compositional formulation.

\paragraph{Fluid momentum.}  From the coefficient of $\delta\mbf{x}_f$,
\begin{equation}
  \rho_f\mbf{a}_f = \rho_f\mbf{g} + \mbf{b}_f - \nabla_{\mbf{x}}(\mu_f J_f) - \lambda\nabla_{\mbf{x}}\phi_f - P^\star_\phi\nabla_{\mbf{x}}\phi_f + \gamma_f\nabla_{\mbf{x}}\phi_f + \phi_f\nabla_{\mbf{x}}\gamma_f + c^+_f\nabla_{\mbf{x}}\tau.
  \label{eq:mom_fluid_raw}
\end{equation}
Using $\mu_f J_f = \phi_f \bar p_f$ and \eqref{eq:pf},
\begin{align*}
  -\nabla_{\mbf{x}}(\phi_f \bar p_f) - \lambda\nabla_{\mbf{x}}\phi_f - P^\star_\phi\nabla_{\mbf{x}}\phi_f + \gamma_f\nabla_{\mbf{x}}\phi_f + \phi_f\nabla_{\mbf{x}}\gamma_f
  &= -\phi_f\nabla_{\mbf{x}}\bar p_f - \bar p_f\nabla_{\mbf{x}}\phi_f \\
  &\quad + \left(-\lambda - P^\star_\phi + \gamma_f\right)\nabla_{\mbf{x}}\phi_f
  + \phi_f\nabla_{\mbf{x}}\gamma_f \\
  &= -\phi_f\nabla_{\mbf{x}}\bar p_f.
\end{align*}
Hence the fluid momentum balance reduces to the phase-pressure form
\begin{equation}
  \boxed{\rho_f\mbf{a}_f = \rho_f\mbf{g} + \mbf{b}_f - \phi_f\nabla_{\mbf{x}}\bar p_f + c^+_f\left(\nabla_{\mbf{x}}\tau - \mbf{v}_f\right),}
  \label{eq:mom_fluid}
\end{equation}
where the reaction-momentum term has been written in the standard Bedford--Drumheller form.

\paragraph{Gas momentum.}  From the coefficient of $\delta\mbf{x}_g$,
\begin{equation}
  \rho_g\mbf{a}_g = \rho_g\mbf{g} + \mbf{b}_g - \nabla_{\mbf{x}}(\mu_g J_g) - \lambda\nabla_{\mbf{x}}\phi_g - P^\star_\phi\nabla_{\mbf{x}}\phi_g + \gamma_g\nabla_{\mbf{x}}\phi_g + \phi_g\nabla_{\mbf{x}}\gamma_g + c^+_g\nabla_{\mbf{x}}\tau.
  \label{eq:mom_gas_raw}
\end{equation}
Using $\mu_g J_g = \phi_g \bar p_g$ and \eqref{eq:pg},
\begin{align*}
  -\nabla_{\mbf{x}}(\phi_g \bar p_g) - \lambda\nabla_{\mbf{x}}\phi_g - P^\star_\phi\nabla_{\mbf{x}}\phi_g + \gamma_g\nabla_{\mbf{x}}\phi_g + \phi_g\nabla_{\mbf{x}}\gamma_g
  &= -\phi_g\nabla_{\mbf{x}}\bar p_g - \bar p_g\nabla_{\mbf{x}}\phi_g \\
  &\quad + \left(-\lambda - P^\star_\phi + \gamma_g\right)\nabla_{\mbf{x}}\phi_g
  + \phi_g\nabla_{\mbf{x}}\gamma_g \\
  &= -\phi_g\nabla_{\mbf{x}}\bar p_g.
\end{align*}
Hence the gas momentum balance reduces to
\begin{equation}
  \boxed{\rho_g\mbf{a}_g = \rho_g\mbf{g} + \mbf{b}_g - \phi_g\nabla_{\mbf{x}}\bar p_g + c^+_g\left(\nabla_{\mbf{x}}\tau - \mbf{v}_g\right).}
  \label{eq:mom_gas}
\end{equation}

\paragraph{Overall momentum.}  Summing the solid, fluid, and gas balances gives the mixture equation.  The direct inter-component body-force pair terms from~\eqref{eq:phi_variation_split} cancel in the sum, but the porosity-conjugate contribution survives through the solid pressure algebra.  Writing the equivalent pore pressure as $\bar p_E \equiv S_f\bar p_f + (1-S_f)\bar p_g-\gamma_{fg}=-\lambda-P^\star_\phi$, the overall balance takes the form
\begin{equation}
  \boxed{\rho\mbf{a} = \nabla_{\mbf{x}}\cdot(\bs\sigma'_s)^T - \nabla_{\mbf{x}}\bar p_E - \nabla_{\mbf{x}}\left[(1-\phi)P^\star_\phi\right] + \rho\mbf{g} - c^+_f\left(\mbf{v}_f - \mbf{v}_s\right) - c^+_g\left(\mbf{v}_g - \mbf{v}_s\right).}
  \label{eq:mom_overall}
\end{equation}
When $\psi_\xi$ does not depend on $\phi$ ($P^\star_\phi = 0$), this reduces to the standard mixture form.

\paragraph{Evolution of $\tau$.}  From the coefficient of $\delta c_\xi$ and choosing to evaluate the functions for the reactant $A$,
\begin{equation}
  \boxed{\left.\p{\tau}{t}\right|_{\mbf{X}_s} = \frac{1}{2}\mbf{v}_s\cdot\mbf{v}_s + \mc{L}_A + \frac{\lambda}{\bar\rho_A}.}
  \label{eq:tau_final}
\end{equation}

\subsection{Connection to the Coleman--Noll procedure}
\label{sec:coleman_connection}

The variational derivation and the Coleman--Noll procedure are complementary.  Hamilton's principle takes the free energies (e.g. $\psi_\xi(\mbf{F}_s, \phi)$) as functions of \emph{independent fields} (subject to constraints) and derives the momentum balances, pressure identifications, and body forces as Euler--Lagrange equations.  The Coleman--Noll procedure (Appendix~\ref{sec:coleman}) starts from the entropy inequality with the \emph{same} energy and derives constitutive restrictions.  The multiplicative decomposition $\mbf{F}_s = \mbf{A}_s \bar{\mbf{F}}_s$ is a constitutive choice imposed within the Coleman--Noll procedure, giving $\psi_\xi(\mbf{F}_s, \phi) = \psi^{\text{iso}}_\xi(\bar{\mbf{F}}_s) + \psi^{\text{dis}}_\xi(J_s, \phi)$, which is a specialization of the general $\psi_\xi(\mbf{F}_s, \phi)$ used in the variational principle.

\section{Coleman-Noll Procedure}
\label{app:coleman_noll}
\label{sec:coleman}

The Coleman-Noll procedure is a systematic thermodynamic approach used to derive restrictions on constitutive models by ensuring compliance with the second law of thermodynamics. This procedure involves formulating the entropy inequality and identifying constraints that must be satisfied by any admissible constitutive relation. We apply this method to our reacting mixture to establish thermodynamically consistent constitutive equations.

Introducing the Helmholtz free energy $\psi_\xi$, the entropy $s_\xi$, and the temperature $\theta_\xi$ for each component $\xi$. These quantities are related to the internal energy through the Legendre transformation
\begin{equation}
  E_\xi = \psi_\xi + s_\xi \theta_{\xi}.
  \label{eq:legendre}
\end{equation}
From \cite{drumheller2000} the conservation of energy for component $\xi$ in the mixture takes the form
\begin{equation}
  \rho_{\xi} \dot{E}_{\xi} + c_{\xi}^{+} E_{\xi} = \mathrm{tr}\left(\left(\bs{\sigma}^{\prime}_{\xi}\right)^{T} \mathbf{L}_{\xi}\right) + \phi_{\xi}\bar{p}_\xi  \frac{\dot{\bar\rho}_{\xi}}{\bar\rho_{\xi}} - \mathcal{L}_{\xi} c_{\xi}^{+} - \nabla_{\mathbf{x}} \cdot \mathbf{q}_{\xi} + \rho_{\xi} r_{\xi} + \dot{e}_{\xi},
  \label{eq:coe}
\end{equation}
where the terms on the right-hand side represent mechanical power, work due to density changes, latent heat of reaction, heat conduction, heat supply, and interaction energy, respectively. $\mathcal{L}_\xi$ is a generalized force that does work when $c^{+}_\xi$ changes.  Substituting the Legendre transformation \eqref{eq:legendre} into the energy equation \eqref{eq:coe} and rearranging to solve for the heat supply term $\rho_\xi r_{\xi}$ yields
\begin{align}
  \rho_{\xi} r_{\xi} =& \, \rho_{\xi} \dot{\psi}_{\xi} + \rho_\xi \dot{s}_\xi \theta_\xi + \rho_\xi s_\xi \dot\theta_\xi + c_{\xi}^{+} \left(\psi_{\xi} + s_\xi \theta_\xi\right) - \notag \\
  & \, \mathrm{tr}(\boldsymbol{\sigma}^{\prime T}_{\xi} \mathbf{L}_{\xi}) - \phi_{\xi} \bar{p}_\xi \frac{\dot{\bar\rho}_{\xi}}{\bar\rho_{\xi}} + \mathcal{L}_{\xi} c_{\xi}^{+} + \nabla_{\mathbf{x}} \cdot \mathbf{q}_{\xi} - \dot{e}_{\xi}, \label{eq:rho_r}
\end{align}
The second law of thermodynamics for the mixture requires that the entropy production rate be non-negative
\begin{equation}
  0 \leq \sum_{\xi} \left[ \rho_{\xi} \dot{s}_{\xi} + \nabla_{\mathbf{x}} \cdot \left( \frac{\mathbf{q}_{\xi}}{\theta_{\xi}} \right) - \frac{\rho_{\xi} r_{\xi}}{\theta_{\xi}} + c_{\xi}^{+} s_{\xi} \right].
  \label{eq:2nd_law}
\end{equation}
To apply the Coleman-Noll procedure, we substitute the expression for the heat supply term \eqref{eq:rho_r} into the entropy inequality \eqref{eq:2nd_law}:
\begin{equation}
  0 \leq \sum_{\xi} \left[ - \frac{1}{\theta_{\xi}} \left( \rho_{\xi} \dot{\psi}_{\xi} +  \rho_\xi s_\xi \dot\theta_\xi + c_{\xi}^{+} \psi_{\xi}  - \mathrm{tr}(\boldsymbol{\sigma}^{\prime T}_{\xi} \mathbf{L}_{\xi}) - \phi_{\xi} \bar{p}_\xi \frac{\dot{\bar\rho}_{\xi}}{\bar\rho_{\xi}} + \mathcal{L}_{\xi} c_{\xi}^{+} - \dot{e}_{\xi}\right) - \frac{\mbf{q}_\xi}{\theta^2_{\xi}} \cdot \nabla_{\mbf{x}} \theta_\xi \right].
  \label{eq:2nd_law_expand_1}
\end{equation}

For our reacting mixture, since $\alpha_s = J_s$ (equation~\eqref{eq:alpha_Js}) implies $\mbf{\bar F}_s = J_s^{-\nicefrac{1}{3}} \mbf{F}_s$ is a known function of $\mbf{F}_s$, we specify the functional dependencies of the Helmholtz free energy as
\begin{align*}
  \psi_A &= \psi_A(\mbf{\bar F}_s(\mbf{F}_s),\, J_s(\mbf{F}_s),\, \phi), \\
  \psi_B &= \psi_B(\mbf{\bar F}_s(\mbf{F}_s),\, J_s(\mbf{F}_s),\, \phi), \\
  \psi_f &= \psi_f(\bar\rho_f), \\
  \psi_g &= \psi_g(\bar\rho_g),
\end{align*}
furthermore, we assume that the solid free energies can be decomposed as
\begin{align*}
  \psi_\xi &= \psi_\xi^{\text{iso}}(\mbf{\bar F}_s(\mbf{F}_s)) + \psi_\xi^{\text{dis}}(J_s(\mbf{F}_s),\, \phi),
\end{align*}
where the first term is the free energy associated with isochoric (i.e.\ volume preserving) deformation of the solid mineral grains, and the second term involves the distention energy, i.e.\ the energy associated with rearranging the incompressible mineral grains during deformation of the solid skeleton.

To focus on the key thermodynamic restrictions and simplify the derivations, we assume a spatially and temporally uniform temperature in what follows.  This allows us to concentrate on the mechanical aspects of the constitutive relationships.

Since $\mbf{\bar F}_s$ and $J_s$ both depend on $\mbf{F}_s$, and $\psi^{\text{dis}}_\xi$ depends on $\phi$, the time derivative of the solid free energies requires the chain rule through all three paths:
\begin{align}
  \dot\psi_\xi = \p{\psi^{\text{iso}}_\xi}{\bar{F}_{mn}} \frac{\partial \bar{F}_{mn}}{\partial F_{kl}} \dot{F}_{kl} + \left.\p{\psi^{\text{dis}}_\xi}{J_s}\right\vert_\phi J_s F^{-T}_{kl} \dot{F}_{kl} + \left.\p{\psi^{\text{dis}}_\xi}{\phi}\right\vert_{J_s}\dot\phi,
  \label{eq:psi_dot_chain}
\end{align}
where
\begin{align}
  \frac{\partial \bar{F}_{mn}}{\partial F_{kl}} = J_s^{-\nicefrac{1}{3}}\left(\delta_{mk}\delta_{nl} - \frac{1}{3} F_{mn} F^{-T}_{kl}\right).
  \label{eq:chain_tensor_app}
\end{align}
The first term in~\eqref{eq:psi_dot_chain} produces the deviatoric (isochoric) contribution to the stress via the projection in~\eqref{eq:chain_tensor_app}.  The second term produces the volumetric (distention) contribution.  The third term, involving $\dot\phi$, will be modified using the mass balance.

For the mechanical power, since $\bs\sigma'_s$ is symmetric (by angular momentum balance), we note
\begin{align}
  \mathrm{tr}\left(\left(\bs\sigma'_s\right)^T \mbf{L}_s\right) &= \sigma'_{ji} \dot{F}_{ik} F^{-1}_{kj} = \frac{1}{J_s} P'_{ik} \dot{F}_{ik},
  \label{eq:mech_power_app}
\end{align}
where $\mbf{P}'_s = J_s \bs\sigma'_s \mbf{F}_s^{-T}$ is the first Piola--Kirchhoff effective stress.

Following \cite{drumheller1980thermomechanical}, the conservation of energy for the entire mixture gives
\begin{align}
  \sum_\xi \dot{e}_\xi &= -\sum_\xi \mbf{b}_\xi \cdot \mbf{v}_\xi \notag \\
  &= -\mbf{b}_f \cdot \left(\mbf{v}_f -  \mbf{v}_s\right) - \mbf{b}_g \cdot \left(\mbf{v}_g - \mbf{v}_s\right).
  \label{eq:disappation}
\end{align}

Substituting~\eqref{eq:psi_dot_chain} and~\eqref{eq:mech_power_app} into the entropy inequality~\eqref{eq:2nd_law_expand_1}, recalling that $\dot{\bar\rho}_A = \dot{\bar\rho}_B = 0$ by incompressibility, we obtain
\begin{align}
  0 \leq & \left(\frac{1}{J_s} P'_{s,ik} - \sum_{\xi=A,B} \rho_\xi \left[\p{\psi^{\text{iso}}_\xi}{\bar{F}_{mn}} \frac{\partial \bar{F}_{mn}}{\partial F_{ik}} + \left.\p{\psi^\text{dis}_\xi}{J_s}\right\vert_\phi J_s F^{-T}_{ik}\right]\right) \dot{F}_{ik} \notag \\
  & + \left(\frac{\phi_f \bar{p}_f}{\bar\rho_f} - \rho_f \p{\psi_f}{\bar\rho_f}\right) \dot{\bar\rho}_f + \left(\frac{\phi_g \bar{p}_g}{\bar\rho_g} - \rho_g \p{\psi_g}{\bar\rho_g}\right) \dot{\bar\rho}_g \notag \\
  & - \sum_{\xi=A,B} \rho_\xi \left.\p{\psi^\text{dis}_\xi}{\phi}\right\vert_{J_s}\dot\phi - \sum_{\xi=A,B,f,g} c^{+}_\xi \left(\mc{L}_\xi + \psi_\xi\right) - \sum_\xi \frac{\dot e_\xi}{\theta}.
  \label{eq:2nd_law_expand_2}
\end{align}
The $\dot\phi$ term is not independently specifiable.  Because the solid constituents are incompressible ($\alpha_s = J_s$), $\phi$ is kinematically determined by $\mbf{F}_s$ and the accumulated reaction.  Specifically,
\begin{equation}
  J_s \dot\phi = (1-\phi)\dot{J}_s - J_s\sum_{\xi=A,B}\frac{c^+_\xi}{\bar\rho_\xi}.
  \label{eq:phi_dot_sub}
\end{equation}
Using $\dot{J}_s = J_s F^{-T}_{ik}\dot{F}_{ik}$ and substituting into~\eqref{eq:2nd_law_expand_2}, the $\dot\phi$ term splits into a contribution proportional to $\dot{F}_{ik}$ (which augments the stress) and a contribution proportional to $c^+_\xi$ (which augments the reaction driving force).  After regrouping:
\begin{align}
  0 \leq & \left(\frac{1}{J_s} P'_{s,ik} - \sum_{\xi=A,B} \rho_\xi \left[\p{\psi^{\text{iso}}_\xi}{\bar{F}_{mn}} \frac{\partial \bar{F}_{mn}}{\partial F_{ik}} + \left(\left.\p{\psi^\text{dis}_\xi}{J_s}\right\vert_\phi + \frac{(1-\phi)}{J_s}\left.\p{\psi^\text{dis}_\xi}{\phi}\right\vert_{J_s}\right) J_s F^{-T}_{ik}\right]\right) \dot{F}_{ik} \notag \\
  & + \left(\frac{\phi_f \bar{p}_f}{\bar\rho_f} - \rho_f \p{\psi_f}{\bar\rho_f}\right) \dot{\bar\rho}_f + \left(\frac{\phi_g \bar{p}_g}{\bar\rho_g} - \rho_g \p{\psi_g}{\bar\rho_g}\right) \dot{\bar\rho}_g \notag \\
  & - \sum_{\xi=A,B,f,g} c^{+}_\xi \left(\mc{L}_\xi + \psi_\xi\right) + \sum_{\xi=A,B}\frac{\rho_\xi}{\bar\rho_\xi} \left.\p{\psi^\text{dis}_\xi}{\phi}\right\vert_{J_s} c^+_\xi - \sum_\xi \frac{\dot e_\xi}{\theta}.
  \label{eq:2nd_law_expand_3}
\end{align}
A sufficient condition for thermodynamic admissibility is that the coefficients of the independent rates $\dot{F}_{ik}$, $\dot{\bar\rho}_f$, $\dot{\bar\rho}_g$, and $\dot{c}_\xi$ each vanish individually.

\paragraph{Coefficient of $\dot{F}_{ik}$:}
\begin{align}
  \frac{1}{J_s} P'_{s,ik} = \sum_{\xi=A,B} \rho_\xi \left[\p{\psi^{\text{iso}}_\xi}{\bar{F}_{mn}} \frac{\partial \bar{F}_{mn}}{\partial F_{ik}} + \left(\left.\p{\psi^{\text{dis}}_\xi}{J_s}\right\vert_\phi + \frac{1-\phi}{J_s}\left.\p{\psi^{\text{dis}}_\xi}{\phi}\right\vert_{J_s}\right) J_s F^{-T}_{ik}\right].
  \label{eq:Fdot_coeff}
\end{align}
Substituting~\eqref{eq:chain_tensor_app}, using $\rho_\xi \partial\psi^{\text{iso}}_\xi / \partial\bar{F}_{mn} = \phi_\xi \bar{P}'_{\xi,mn}$ where $\mbf{\bar P}'_\xi = \partial W^{\text{iso}}_\xi / \partial\mbf{\bar F}_s$ is the undistended intrinsic first Piola--Kirchhoff effective stress, and defining the component distention pressure
\begin{align}
  \bar{P}^\star_\xi &= \left.\rho_\xi \frac{\partial \psi^{\text{dis}}_\xi}{\partial J_s}\right\vert_\phi + \frac{1-\phi}{J_s}\left.\rho_\xi\frac{\partial \psi^{\text{dis}}_\xi}{\partial\phi}\right\vert_{J_s}, \\
   &= \left.\frac{\partial W^{\text{dis}}_\xi}{\partial J_s}\right\vert_\phi + \frac{1-\phi}{J_s}\left.\frac{\partial W^{\text{dis}}_\xi}{\partial\phi}\right\vert_{J_s},
\end{align}
this gives
\begin{align}
  \mbf{P}'_s = \sum_{\xi=A,B} \left[J_s^{\nicefrac{2}{3}} \phi_\xi \left(\mbf{\bar P}'_\xi - \frac{1}{3}\left(\mbf{\bar P}'_\xi : \mbf{F}_s\right) \mbf{F}_s^{-T}\right) + \phi_\xi \bar{P}^\star_\xi J_s\, \mbf{F}_s^{-T}\right].
  \label{eq:P_CN_direct}
\end{align}
The corresponding Cauchy effective stress is
\begin{align}
  \bs\sigma'_s = \sum_{\xi=A,B}\left[J_s^{-\nicefrac{1}{3}} \phi_\xi \left(\mbf{\bar P}'_\xi \mbf{F}_s^T - \frac{1}{3}\left(\mbf{\bar P}'_\xi : \mbf{F}_s\right) \mbf{I}\right) + \phi_\xi \bar{P}^\star_\xi \mbf{I}\right].
  \label{eq:sigma_CN}
\end{align}

\paragraph{Coefficient of $\dot{\bar\rho}_f$:}
\begin{align}
  \bar{p}_f = \bar\rho_f^2 \p{\psi_f}{\bar\rho_f}.
  \label{eq:thermo_pressure}
\end{align}

\paragraph{Coefficient of $\dot{\bar\rho}_g$:}
\begin{align}
  \bar{p}_g = \bar\rho_g^2 \p{\psi_g}{\bar\rho_g}.
  \label{eq:thermo_pressure_g}
\end{align}

\paragraph{Residual inequality (reactions).}  If the constitutive restrictions~\eqref{eq:P_CN_direct}, \eqref{eq:thermo_pressure}, and \eqref{eq:thermo_pressure_g} are satisfied, the sufficient condition for thermodynamic admissibility away from equilibrium ($\dot{c}_\xi \ne 0$) reduces to
\begin{align}
  0 \leq \sum_{\xi = A, B} c^{+}_\xi \left(\phi_\xi \left.\p{\psi^\text{dis}_\xi}{\phi}\right\vert_{J_s} - \mc{L}_\xi - \psi_\xi\right) - c^{+}_f \left(\mc{L}_f + \psi_f\right) - c_g^+ \left(\mc{L}_g + \psi_g\right).
  \label{eq:thermo_reaction}
\end{align}
This leads us to our choice for $\mathcal{L}_A$ in the $\tau$ evolution equation as
\begin{align}
 \mc{L}_A = \phi_A \left.\p{\psi^\text{dis}_A}{\phi}\right\vert_{J_s} - \psi_A.
\end{align}

\paragraph{Evaluation for $W^{\text{dis}}_\xi(J_s, \phi) = K^\star_\xi (\ln J_s)^2 / [2(1-\phi)]$.}  The required partial derivatives are
\begin{align}
  \left.\p{W^{\text{dis}}_\xi}{J_s}\right\vert_\phi &= \frac{K^\star_\xi \ln J_s}{(1-\phi)\,J_s}, &
  \left.\p{W^{\text{dis}}_\xi}{\phi}\right\vert_{J_s} &= \frac{K^\star_\xi (\ln J_s)^2}{2(1-\phi)^2}.
  \label{eq:Wdis_partials}
\end{align}
The partial distention pressure is therefore
\begin{equation}
  \bar{P}^\star_\xi = \frac{K^\star_\xi \ln J_s}{(1-\phi)\,J_s}\left(1 + \frac{1}{2}\ln J_s\right).
  \label{eq:Pstar_explicit}
\end{equation}
In the Cauchy effective stress, the effective bulk modulus from the mixture rule gives
\begin{equation}
  \bs\sigma'_s = G_{\text{eff}}\, J_s^{-\nicefrac{2}{3}} \mbox{dev}(\mbf{B}_s) + \frac{K^\star_{\text{eff}} \ln J_s}{J_s}\left(1 + \frac{\ln J_s}{2}\right) \mbf{I}.
  \label{eq:sigma_CN_explicit}
\end{equation}
The porosity-conjugate pressure appearing in the variational principle is
\begin{equation}
  P^\star_\phi = \sum_{\xi=A,B} \rho_\xi \left.\p{\psi^{\text{dis}}_\xi}{\phi}\right\vert_{J_s} = \frac{K^\star_{\text{eff}} (\ln J_s)^2}{2(1-\phi)}.
  \label{eq:Pstar_phi_explicit}
\end{equation}

%\input{appendix_no_mechanics}

\bibliography{all}
\bibliographystyle{plain}
\end{document}
